\documentclass[11pt]{amsart}
\usepackage[margin=1.1in]{geometry}
\usepackage{amsmath,amssymb,amsthm}
\usepackage[hidelinks]{hyperref}
\usepackage{color}
\usepackage{fancyhdr}
% Banner on every page (including the first).
\newcommand{\aibanner}{\centerline{\textcolor{red}{\bfseries This paper is entirely AI-generated}}}
\pagestyle{fancy}
\fancyhf{}
\fancyhead[C]{\aibanner}
\fancyfoot[C]{\thepage}
\renewcommand{\headrulewidth}{0pt}
\fancypagestyle{plain}{\fancyhf{}\fancyhead[C]{\aibanner}\fancyfoot[C]{\thepage}\renewcommand{\headrulewidth}{0pt}}
\fancypagestyle{firstpage}{\fancyhf{}\fancyhead[C]{\aibanner}\fancyfoot[C]{\thepage}\renewcommand{\headrulewidth}{0pt}}
\setlength{\headheight}{14pt}

\theoremstyle{plain}
\newtheorem{theorem}{Theorem}[section]
\newtheorem{proposition}[theorem]{Proposition}
\newtheorem{lemma}[theorem]{Lemma}
\newtheorem{corollary}[theorem]{Corollary}
\theoremstyle{definition}
\newtheorem{definition}[theorem]{Definition}
\newtheorem{construction}[theorem]{Construction}
\theoremstyle{remark}
\newtheorem{remark}[theorem]{Remark}
\newtheorem{facts}[theorem]{Facts}

\newcommand{\N}{\mathbb N}
\newcommand{\Q}{\mathbb Q}
\newcommand{\M}{\mathrm M}
\newcommand{\free}[1]{k\langle #1\rangle}
\newcommand{\ff}[1]{\mathcal F(#1)}
\newcommand{\Hom}{\operatorname{Hom}}
\newcommand{\op}{\mathrm{op}}
\newcommand{\id}{\mathrm{id}}
\newcommand{\Fix}{\operatorname{Fix}}
\newcommand{\Rt}{\widetilde R}
\newcommand{\Kt}{\widetilde K}
\newcommand{\st}{\tilde\sigma}
\newcommand{\dt}{\tilde\delta}
\newcommand{\T}{\mathcal T}

\title{PCI rings and V-domains are not left-right symmetric}
\date{September 2026}

\begin{document}

\begin{abstract}
We construct a countable domain $R$ which is a left PCI domain (every proper cyclic left $R$-module is injective; in particular $R$ is a simple principal left ideal domain and a left V-domain) but which is not right Ore. Consequently $R$ is neither a right V-ring nor a right PCI ring, and its opposite ring $R^{\op}$ is a right PCI domain which is not left PCI, not left V, and not left Ore. This answers, in the negative, questions of Cozzens and Faith, of Damiano, and of Jain, Lam and Leroy. The ring is an Ore extension $\Kt[t;\st,\dt]$ over a division ring $\Kt$ with a non-surjective endomorphism $\st$; the division ring is a free field in the sense of Cohn, chosen so that every linear equation of the kind that governs divisibility of the cyclic modules $R/Rf$ has a solution.
\end{abstract}

\maketitle

\section{Introduction}

A ring $R$ is a \emph{right PCI ring} if every proper cyclic right $R$-module, that is, every cyclic right module not isomorphic to $R_R$, is injective; it is a \emph{right V-ring} if every simple right $R$-module is injective. Faith \cite[Theorem]{Fai73} showed that a right PCI ring is either semisimple Artinian or a simple right semihereditary right Ore domain, and Damiano \cite[Theorem~1]{Dam79} showed that it is then right Noetherian. Examples of (two-sided) PCI domains were given by Cozzens \cite{Coz70}. Cozzens and Faith \cite[6.25, 6.26]{CF75} (originally \cite[Theorem~22]{Fai73}) showed, using a theorem of Boyle \cite{Boy74}, that a right Noetherian right PCI domain which is \emph{left Ore} is left Noetherian and then also left PCI; by Damiano's theorem the hypothesis that it be right Noetherian is automatic. Whether every right PCI ring is left PCI has been open since \cite{Fai73} (see also \cite[\S2]{BG75} and \cite{CF75}), and is restated as open in \cite[\S2]{JS09}: the question is equivalent to asking whether a right PCI domain must be left Ore \cite{Dam79}, and its cousin, whether a left V-domain must be a right V-domain, goes back to Cozzens and Faith \cite[\S7, (10)]{CF75} and is restated as open in \cite[\S1]{JLL09} and \cite[Remark~2.8]{BDV18}. Cozzens and Faith asked for ``a right V-domain which is not a left Ore domain'' \cite[\S7, (7)]{CF75} and for ``a right Noetherian right V-(PCI-)domain which is not a left V-(PCI-)domain'' \cite[\S7, (10)]{CF75}, and proposed to look for one among Ore extensions $K[t;\sigma,\delta]$ over a division ring $K$ with a non-surjective endomorphism $\sigma$, an idea taken up in \cite[\S6]{JLL09}, where it is shown that for a left V-domain of the form $K[t;\sigma,\delta]$ the ring is right PCI if and only if $\sigma$ is onto, and where the authors state that they have no example of a left V-domain $K[t;\sigma,\delta]$ with $\sigma$ not onto.

We produce such an example.

\begin{theorem}[Main Theorem]\label{thm:main}
Let $k$ be a countable commutative field. There is a countable division ring $\Kt\supseteq k$, an injective but non-surjective endomorphism $\st$ of $\Kt$ and a $\st$-derivation $\dt$ of $\Kt$ such that the Ore extension $\Rt=\Kt[t;\st,\dt]$ has the following properties.
\begin{enumerate}
\item $\Rt$ is a simple principal left ideal domain and a left PCI domain: every proper cyclic left $\Rt$-module is injective. In particular $\Rt$ is a left V-domain.
\item $\Rt$ is not right Ore. Consequently $\Rt$ is not a right PCI ring, and not a right V-ring.
\end{enumerate}
Hence $\Rt^{\op}$ is a simple right Noetherian right hereditary right PCI domain which is not left Ore, not a left V-ring, and not left PCI.
\end{theorem}

The strategy is the one suggested in \cite[\S7, (10)]{CF75}: divisibility of the cyclic modules $R/Rf$ over $R=K[t;\sigma,\delta]$ is governed by linear equations in $K^n$ involving $\sigma$ and $\delta$ (Section~\ref{sec:ore}), and one has to enlarge $K$ until all such equations have solutions while keeping $\sigma$ non-surjective. The enlargement is done with Cohn's free fields (Section~\ref{sec:freefield}): the solutions are adjoined as free generators, which guarantees that nothing collapses, and the resulting free field carries an extension of $\sigma$ and $\delta$ because both are prescribed on generators. This is the same manoeuvre that Cohn uses in \cite[\S5]{Coh77} and \cite[\S5.9]{CohSF}. Non-surjectivity of $\sigma$ in the limit is checked with a single automorphism of the free field (Lemma~\ref{lem:nonsurj}). The only results quoted without proof are the standard ones on Ore extensions and PCI rings from \cite{CF75,Fai73,JLL09} and the following facts from Cohn's theory: a free algebra is a fir, a semifir has a universal field of fractions, and honest homomorphisms extend to universal fields of fractions.

\subsection*{Conventions}
Rings have an identity, and homomorphisms preserve it. A \emph{division ring} is a ring in which every non-zero element is invertible. Throughout, $k$ is a commutative field. For a division ring $K$, an \emph{endomorphism} $\sigma$ of $K$ is a ring endomorphism (automatically injective), and a \emph{$\sigma$-derivation} of $K$ is an additive map $\delta\colon K\to K$ with
\[
\delta(ab)=\sigma(a)\delta(b)+\delta(a)b\qquad(a,b\in K).
\]
For $u\ne0$ one has $\delta(u^{-1})=-\sigma(u)^{-1}\delta(u)u^{-1}$, since $0=\delta(uu^{-1})=\sigma(u)\delta(u^{-1})+\delta(u)u^{-1}$. A subset $D$ of a division ring is a \emph{division subring} if it is a subring closed under inverses of non-zero elements; the division subring \emph{generated} by a subset $S$ is the smallest division subring containing $S$. Two facts will be used constantly: the equalizer $\{u:\alpha(u)=\beta(u)\}$ of two ring homomorphisms $\alpha,\beta$ out of a division ring is a division subring; and if $\alpha=\beta$, the equalizer $\{u:\delta(u)=\delta'(u)\}$ of an $\alpha$-derivation $\delta$ and an $\alpha$-derivation $\delta'$ (both defined on the domain of $\alpha$ with values in the codomain, satisfying the Leibniz rule with respect to $\alpha$) is a division subring, by the formula for $\delta(u^{-1})$.

\section{Ore extensions over division rings}\label{sec:ore}

Let $K$ be a division ring with an endomorphism $\sigma$ and a $\sigma$-derivation $\delta$. The \emph{Ore extension} $R=K[t;\sigma,\delta]$ is the ring which is a free left $K$-module with basis $1,t,t^2,\dots$ and multiplication determined by
\[
ta=\sigma(a)t+\delta(a)\qquad(a\in K).
\]
Elements are written $f=\sum_i f_it^i$ with coefficients on the left, and $\deg f$ is the usual degree. Since $\sigma$ is injective, $\deg(fg)=\deg f+\deg g$, so $R$ is a domain, and the units of $R$ are the non-zero elements of $K$. For $f\ne0$ and $r\in R$ there are unique $q,b\in R$ with $r=qf+b$ and $\deg b<\deg f$ (left division algorithm; compare \cite[Theorem~1.3.2]{CohFIR}, where coefficients are written on the right, so that all statements there are the mirror images of the ones here). Hence every left ideal of $R$ is principal, $R$ is a \emph{principal left ideal domain}, and $R$ is \emph{atomic}: every non-zero non-unit is a finite product of atoms (irreducible elements), by induction on the degree.

\subsection{The modules $R/Rf$}
Let $f=t^n+f_{n-1}t^{n-1}+\dots+f_0$ be monic of degree $n\ge1$. Every element of $R$ is uniquely of the form $qf+b$ with $\deg b<n$, so the classes of $1,t,\dots,t^{n-1}$ form a basis of $R/Rf$ as a left $K$-vector space, and we identify $R/Rf$ with the space $K^n$ of row vectors by
\[
b=\sum_{i=0}^{n-1}b_it^i+Rf\ \longmapsto\ \bar b=(b_0,\dots,b_{n-1}).
\]
Let $C_f\in\M_n(K)$ be the companion matrix of $f$, with rows $e_2,\dots,e_n,(-f_0,\dots,-f_{n-1})$, where $e_1,\dots,e_n$ are the standard basis vectors. For a row vector $v$ over any ring containing $K$ to which $\sigma$ and $\delta$ have been extended, we write $\sigma(v)$ and $\delta(v)$ for the entrywise images, and define
\[
T_f(v)=\sigma(v)C_f+\delta(v).
\]
The map $T_f$ is additive and \emph{pseudo-linear}: $T_f(av)=\sigma(a)T_f(v)+\delta(a)v$ for $a\in K$. For $g=\sum_{i=0}^m g_it^i\in R$ we put $g(T_f)(v)=\sum_{i=0}^m g_i\,T_f^i(v)$, where $T_f^i$ is the $i$-fold iterate and $g_i$ acts by scalar multiplication.

\begin{lemma}\label{lem:action}
Let $f$ be monic of degree $n\ge1$, $g\in R$ and $b\in R$ with $\deg b<n$. Then the coefficient vector of the remainder of $gb$ modulo $Rf$ is $g(T_f)(\bar b)$. In particular
\[
gb\equiv 1\pmod{Rf}\iff g(T_f)(\bar b)=e_1 .
\]
\end{lemma}

\begin{proof}
Left multiplication by $a\in K$ on $R/Rf$ corresponds to scalar multiplication of $\bar b$ by $a$. Left multiplication by $t$ sends $\sum_{i<n}b_it^i$ to $\sum_{i<n}\sigma(b_i)t^{i+1}+\sum_{i<n}\delta(b_i)t^i$, and $t^n\equiv-\sum_{j<n}f_jt^j\pmod{Rf}$; so the coefficient vector of $tb+Rf$ is $\sigma(\bar b)C_f+\delta(\bar b)=T_f(\bar b)$. Iterating and adding gives the claim for $g=\sum g_it^i$; the last statement is the case where the remainder is $1$, whose coefficient vector is $e_1$.
\end{proof}

This is the identification of $R/Rf$ with $(K^n,T_f)$ used in \cite[\S4]{JLL09}.

\subsection{The criterion}
Recall that a domain $R$ is a \emph{left PCI domain} if every proper cyclic left $R$-module is injective, and a \emph{left V-domain} if every simple left $R$-module is injective.

\begin{proposition}\label{prop:criterion}
Let $R=K[t;\sigma,\delta]$ and suppose that for every monic $f\in R$ of degree $n\ge1$ and every $g\in R$ of degree $\ge1$ there is $v\in K^n$ with $g(T_f)(v)=e_1$. Then $R=Rf+gR$ for all non-zero $f,g\in R$. Consequently $R$ is a simple ring, a left PCI domain and a left V-domain.
\end{proposition}

\begin{proof}
Let $f,g\ne0$. If $\deg g=0$ then $g$ is a unit and $gR=R$; if $\deg f=0$ then $Rf=R$. Otherwise write $f=af_1$ with $a\in K^\times$ and $f_1$ monic, so that $Rf=Rf_1$, choose $v$ with $g(T_{f_1})(v)=e_1$ and put $b=\sum_{i<n}v_it^i$. By Lemma~\ref{lem:action}, $gb-1\in Rf_1=Rf$, so $1\in Rf+gR$ and $R=Rf+gR$.

If $I$ is a non-zero two-sided ideal and $0\ne f\in I$, then $R=Rf+fR\subseteq I$; so $R$ is simple. Since $R$ is an atomic principal left ideal domain, the equivalence of ``$R=Rf+gR$ for all non-zero $f,g$'' with ``$R$ is a left PCI domain'' and with ``$R$ is a left V-domain'' is \cite[Theorem~3.2, (iv)$\Leftrightarrow$(i)$\Leftrightarrow$(ii)]{JLL09}.
\end{proof}

For completeness we indicate why $R=Rf+gR$ for all $f,g$ forces every $R/Rf$ to be injective; this is the content of the cited theorem. Over the principal left ideal domain $R$, a left module $E$ is injective if and only if it is divisible, by Baer's criterion; so $R/Rf$ is injective if and only if for every $g\ne0$ and every $c\in R$ there is $b$ with $gb\equiv c\pmod{Rf}$. When $f$ is an atom, $R/Rf$ is simple, and for $c\notin Rf$ the left ideal $\{r:rc\in Rf\}$ is $Rf'$ for an atom $f'$; from $gb'\equiv1\pmod{Rf'}$ one gets $gb'c\equiv c\pmod{Rf}$. For a general $f=f_1f_2$, the module $R/Rf$ is an extension of $R/Rf_2$ by $Rf_2/Rf_1f_2\cong R/Rf_1$, and an extension of injective modules is injective; induction on the number of atomic factors finishes the argument. The proper cyclic left modules are exactly the $R/Rf$ with $f\ne0$, since $R/Rf$ is finite-dimensional over $K$ while $R$ is not.

\subsection{The right side}

\begin{proposition}\label{prop:notore}
Let $R=K[t;\sigma,\delta]$ with $\sigma$ not surjective, and let $c\in K\setminus\sigma(K)$. Then $tR\cap ctR=0$. In particular $R$ is not right Ore, and $R$ is not a right PCI ring.
\end{proposition}

\begin{proof}
Suppose $tr=ctr'$ with $r,r'\ne0$; comparing degrees, $\deg r=\deg r'=p$, say, with leading coefficients $r_p,r'_p$. The leading coefficient of $tr$ is $\sigma(r_p)$ and that of $ctr'$ is $c\sigma(r'_p)$, so $c=\sigma(r_pr_p'^{-1})\in\sigma(K)$, a contradiction. Thus $tR\cap ctR=0$ and $R$ is not right Ore. A right PCI domain is right Ore by \cite[6.17]{CF75} (see also \cite[Theorem]{Fai73}), and $R$ is a domain; so $R$ is not right PCI.
\end{proof}

This is \cite[Proposition~6.1]{JLL09} in the form we need. For the failure of the right V-property we use a cardinality argument in the spirit of Lawrence \cite[Theorem~4]{Law77}.

\begin{lemma}\label{lem:countable}
Let $R$ be a countable domain containing non-zero elements $a,b$ with $aR\cap bR=0$. Then no non-zero countable right $R$-module is injective. In particular no simple right $R$-module is injective, and $R$ is not a right V-ring.
\end{lemma}

\begin{proof}
The right ideals $b^iaR$ ($i\in\N$) are independent. Indeed, if $\sum_{i=0}^N b^iar_i=0$, then
\[
ar_0=-b\sum_{i\ge1}b^{i-1}ar_i\ \in\ aR\cap bR=0,
\]
so $r_0=0$ because $R$ is a domain; cancelling $b$ and repeating gives $r_i=0$ for all $i$. Each $b^iaR\cong R_R$, so $I=\bigoplus_{i\in\N}b^iaR$ is a free right ideal of countably infinite rank. If $E$ is an injective right $R$-module, the restriction map $\Hom_R(R,E)\to\Hom_R(I,E)$ is surjective; since $\Hom_R(R,E)\cong E$ and $\Hom_R(I,E)\cong E^{\N}$, this gives $|E^{\N}|\le|E|$. For a non-zero countable $E$ this is impossible, since $|E^{\N}|\ge2^{\aleph_0}>\aleph_0\ge|E|$. Finally, a simple right module over a countable ring is a cyclic module, hence countable.
\end{proof}

\subsection{Reduction}
By Propositions~\ref{prop:criterion} (applied with $e=e_1$) and~\ref{prop:notore} and Lemma~\ref{lem:countable}, the Main Theorem follows from:

\begin{theorem}\label{thm:target}
Let $k$ be a commutative field. There is a division ring $\Kt\supseteq k$ with an endomorphism $\st$ which is not surjective and a $\st$-derivation $\dt$, such that for every monic $f\in\Kt[t;\st,\dt]$ of degree $n\ge1$, every $g\in\Kt[t;\st,\dt]$ of degree $\ge1$ and every $e\in\Kt^n$ there is $v\in\Kt^n$ with $g(T_f)(v)=e$. If $k$ is countable, $\Kt$ can be taken countable.
\end{theorem}

For the right-hand statements of the Main Theorem we shall also use the following observation, which lets Lemma~\ref{lem:countable} do the work of \cite[6.17]{CF75} in the countable case.

\begin{lemma}\label{lem:pciv}
Let $R$ be a domain which is not a division ring. If $R$ is a right PCI ring then $R$ is a right V-ring; if $R$ is a left PCI ring then $R$ is a left V-ring.
\end{lemma}

\begin{proof}
A simple right $R$-module is $R/M$ for a maximal right ideal $M$. If $R/M\cong R_R$, then the composite $R\to R/M\cong R$ is a surjective endomorphism of $R_R$, hence left multiplication by an element $c$ with $cd=1$ for some $d$; its kernel $\{x:cx=0\}$ is $0$ because $R$ is a domain and $c\ne0$. So $M=0$, and $R_R$ is simple: every non-zero $x$ satisfies $xR=R$, so $xy=1$ for some $y$, and then $(yx-1)y=y(xy)-y=0$ gives $yx=1$ as $y\ne0$. Thus $R$ would be a division ring. Hence $R/M$ is a proper cyclic module, and it is injective when $R$ is right PCI. The left-handed statement is the mirror image.
\end{proof}

The rest of the paper proves Theorem~\ref{thm:target}.

\section{Free fields}\label{sec:freefield}

We recall what is needed from Cohn's theory of universal fields of fractions \cite[Chapter~4 and \S5.4]{CohSF}, \cite[Chapters~2, 7]{CohFIR}. A square matrix $A$ of size $p$ over a ring $S$ is \emph{full} if it cannot be written as $A=PQ$ with $P$ of size $p\times r$, $Q$ of size $r\times p$ and $r<p$. A homomorphism $\phi\colon S\to A$ of rings is \emph{fully inverting} if $\phi(A)$ is invertible over $A$ for every full matrix $A$ over $S$; here $\phi(A)$ is the matrix obtained by applying $\phi$ entrywise. A homomorphism $S\to S'$ between rings is \emph{honest} if it maps full matrices over $S$ to full matrices over $S'$. For a set $Y$, $\free Y$ denotes the free associative $k$-algebra on $Y$.

\begin{facts}\label{facts}
Let $S$ be a semifir (for instance a fir).
\begin{enumerate}
\item[(F1)] The free algebra $\free Y$ is a fir, for every set $Y$ \cite[Corollary~2.5.2]{CohFIR}, \cite[Theorem~5.4.1]{CohSF}.
\item[(F2)] $S$ has a \emph{universal field of fractions}: a division ring $\mathcal U(S)$ with a fully inverting homomorphism $\lambda\colon S\to\mathcal U(S)$, such that every fully inverting homomorphism $\phi\colon S\to A$ into an arbitrary ring $A$ factors uniquely as $\phi=\phi'\circ\lambda$ with $\phi'\colon\mathcal U(S)\to A$ a ring homomorphism. Indeed $\mathcal U(S)$ is the universal localization of $S$ at the set of all full matrices, which is a division ring \cite[Theorem~4.5.8, Corollary~4.5.9]{CohSF}, \cite[Theorem~7.5.13]{CohFIR}; the factorization property is the defining property of a universal localization \cite[\S4.1]{CohSF}. The map $\lambda$ is injective, because a $1\times1$ matrix $(a)$ with $a\ne0$ is full over a domain.
\item[(F3)] If $S,S'$ are semifirs and $\phi\colon S\to S'$ is honest, then $\phi$ extends uniquely to a homomorphism $\mathcal U(S)\to\mathcal U(S')$ \cite[Theorem~4.5.10]{CohSF}. The extension is injective, being a homomorphism between division rings.
\end{enumerate}
\end{facts}

We identify $S$ with its image $\lambda(S)$ in $\mathcal U(S)$. The following consequence of (F2) will be used repeatedly.

\begin{lemma}\label{lem:generated}
$\mathcal U(S)$ is generated by $S$ as a division ring. Consequently, if $S$ is countable, so is $\mathcal U(S)$.
\end{lemma}

\begin{proof}
Let $D\subseteq\mathcal U(S)$ be the division subring generated by $S$. A square matrix over $D$ which is invertible over $\mathcal U(S)$ is invertible over $D$ (its rows are left linearly independent over $\mathcal U(S)$, hence over $D$, so it is invertible over the division ring $D$). Thus the corestriction $S\to D$ is fully inverting and factors through $\lambda$ as $\psi\colon\mathcal U(S)\to D$. The composite of $\psi$ with the inclusion $D\subseteq\mathcal U(S)$ is a homomorphism $\mathcal U(S)\to\mathcal U(S)$ which is the identity on $S$, hence equals the identity by the uniqueness in (F2). So $D=\mathcal U(S)$. The division subring generated by a countable set is countable, being the union of the countably many sets obtained by iterating the ring operations and inversion.
\end{proof}

\begin{lemma}\label{lem:retract}
Let $\phi\colon S\to S'$ and $\rho\colon S'\to S$ be ring homomorphisms with $\rho\circ\phi=\id_S$. Then $\phi$ is honest.
\end{lemma}

\begin{proof}
If $A$ is full over $S$ and $\phi(A)=PQ$ with $r<p$, then $A=\rho\phi(A)=\rho(P)\rho(Q)$ is not full. This is \cite[Proposition~4.5.1]{CohSF}.
\end{proof}

\begin{definition}
For a set $Y$ we write $\ff Y=\mathcal U(\free Y)$ and call it the \emph{free field} on $Y$ over $k$. By (F1), (F2) it exists, contains $\free Y$, and is generated by $Y$ as a division ring (Lemma~\ref{lem:generated}). The field $k$ is central in $\ff Y$: the centralizer of $k$ is a division subring containing $\free Y$.
\end{definition}

\begin{lemma}\label{lem:subset}
Let $Y'\subseteq Y$. The inclusion $\free{Y'}\to\free Y$ extends uniquely to an embedding $\ff{Y'}\to\ff Y$, whose image is the division subring of $\ff Y$ generated by $Y'$. If $Y=\bigcup_iY_i$ is a directed union, then, after these identifications, $\ff Y=\bigcup_i\ff{Y_i}$.
\end{lemma}

\begin{proof}
The inclusion has the retraction $\free Y\to\free{Y'}$ which sends the elements of $Y\setminus Y'$ to $0$, so it is honest by Lemma~\ref{lem:retract} and extends to an embedding by (F3). Its image is a division subring containing $Y'$, and it is generated by $Y'$ because $\ff{Y'}$ is generated by $Y'$ (Lemma~\ref{lem:generated}). For the last claim, $\bigcup_i\ff{Y_i}$ is a division subring of $\ff Y$ (the union is directed) containing $Y$, hence is all of $\ff Y$.
\end{proof}

Henceforth we regard $\ff{Y'}$ as a division subring of $\ff Y$ whenever $Y'\subseteq Y$.

\begin{proposition}[endomorphisms and derivations prescribed on generators]\label{prop:extend}
Let $Y$ be a set, $s\colon Y\to Y$ an injective map and $d\colon Y\to\ff Y$ an arbitrary map. Then there is exactly one pair $(\sigma,\delta)$ consisting of an endomorphism $\sigma$ of $\ff Y$ and a $\sigma$-derivation $\delta$ of $\ff Y$ such that
\[
\sigma(y)=s(y),\qquad\delta(y)=d(y)\qquad(y\in Y).
\]
Moreover, if $Y'\subseteq Y$ satisfies $s(Y')\subseteq Y'$ and $d(Y')\subseteq\ff{Y'}$, then $\sigma(\ff{Y'})\subseteq\ff{Y'}$ and $\delta(\ff{Y'})\subseteq\ff{Y'}$, and the restrictions of $\sigma,\delta$ to $\ff{Y'}$ are the pair determined by $(Y',s|_{Y'},d|_{Y'})$.
\end{proposition}

\begin{proof}
Write $U=\ff Y$ and $S=\free Y$.

\emph{Step 1: the endomorphism $\hat s$ of $S$ induced by $s$ is honest.} $\hat s$ maps $S$ isomorphically onto the subalgebra $\free{s(Y)}$, and $\free{s(Y)}\subseteq S$ has the retraction $r\colon S\to\free{s(Y)}$ sending the generators outside $s(Y)$ to $0$. Then $\rho=\hat s^{-1}\circ r$ (with $\hat s^{-1}\colon\free{s(Y)}\to S$ the inverse of the isomorphism) satisfies $\rho\circ\hat s=\id_S$, so $\hat s$ is honest by Lemma~\ref{lem:retract}. Consequently, for every full matrix $A$ over $S$, the matrix $\hat s(A)$ is full, hence invertible over $U$ by (F2).

\emph{Step 2: a homomorphism into $\M_2(U)$.} Since $S$ is free and $k$ is central in $U$, there is a unique $k$-algebra homomorphism $\theta\colon S\to\M_2(U)$ with
\[
\theta(y)=\begin{pmatrix} s(y)& d(y)\\ 0& y\end{pmatrix}\qquad(y\in Y).
\]
The set of $a\in S$ for which $\theta(a)$ is upper triangular with diagonal $(\hat s(a),a)$ is a $k$-subalgebra containing $Y$, so $\theta(a)=\begin{pmatrix}\hat s(a)&D(a)\\0&a\end{pmatrix}$ for all $a\in S$ and suitable $D(a)\in U$. For a $p\times p$ matrix $A$ over $S$, the matrix $\theta(A)\in\M_p(\M_2(U))$ is, after the obvious permutation of rows and columns, the block matrix $\begin{pmatrix}\hat s(A)&D(A)\\0&A\end{pmatrix}\in\M_{2p}(U)$. If $A$ is full then both diagonal blocks are invertible over $U$ (Step 1 and (F2)), hence so is $\theta(A)$. Thus $\theta$ is fully inverting and by (F2) factors uniquely as $\theta=\theta_U\circ\lambda$ with a ring homomorphism $\theta_U\colon U\to\M_2(U)$.

\emph{Step 3: reading off $\sigma$ and $\delta$.} Let $V$ be the set of $u\in U$ for which $\theta_U(u)$ is upper triangular with $(2,2)$-entry $u$. Products and sums of such matrices are again of this form, so $V$ is a subring of $U$, and $V\supseteq S$. If $0\ne u\in V$, say $\theta_U(u)=\begin{pmatrix}a&b\\0&u\end{pmatrix}$, then this matrix is invertible in $\M_2(U)$ (its inverse is $\theta_U(u^{-1})$), which forces $a\ne0$ (otherwise the row vector $(1,-bu^{-1})$ would annihilate it from the left); and then $\theta_U(u^{-1})=\begin{pmatrix}a^{-1}&-a^{-1}bu^{-1}\\0&u^{-1}\end{pmatrix}$, so $u^{-1}\in V$. Hence $V$ is a division subring containing $S$, and $V=U$ by Lemma~\ref{lem:generated}. Write
\[
\theta_U(u)=\begin{pmatrix}\sigma(u)&\delta(u)\\0&u\end{pmatrix}\qquad(u\in U).
\]
That $\theta_U$ is a ring homomorphism says exactly that $\sigma$ is a ring endomorphism of $U$ and $\delta$ is a $\sigma$-derivation (\cite[\S2.1, (8)--(9)]{CohSF}; the same device is used in \cite[Theorem~7.5.17]{CohFIR} for ordinary derivations of Sylvester domains). By construction $\sigma(y)=s(y)$ and $\delta(y)=d(y)$ for $y\in Y$.

\emph{Uniqueness.} If $(\sigma',\delta')$ is another such pair, the equalizer of $\sigma$ and $\sigma'$ is a division subring of $U$ containing $Y$, hence equals $U$; and then the equalizer of $\delta$ and $\delta'$ is a division subring containing $Y$, hence equals $U$.

\emph{Restriction.} Let $U'=\ff{Y'}\subseteq U$ and $W=\{u\in U:\sigma(u)\in U',\ \delta(u)\in U'\}$. By the Leibniz rule and the formula for $\delta(u^{-1})$, $W$ is a division subring of $U$; it contains $Y'$ because $s(Y')\subseteq Y'\subseteq U'$ and $d(Y')\subseteq U'$. Hence $W\supseteq U'$, so $U'$ is stable under $\sigma$ and $\delta$. The restrictions are an endomorphism and a $\sigma$-derivation of $U'$ with the prescribed values on $Y'$, so they coincide with the pair determined by $(Y',s|_{Y'},d|_{Y'})$ by the uniqueness just proved, applied to $Y'$.
\end{proof}

\begin{lemma}[non-surjectivity]\label{lem:nonsurj}
In the situation of Proposition~\ref{prop:extend}, let $y_0\in Y\setminus s(Y)$. Then $y_0\notin\sigma(\ff Y)$.
\end{lemma}

\begin{proof}
Let $\tau_0$ be the $k$-algebra automorphism of $\free Y$ with $\tau_0(y_0)=y_0+1$ and $\tau_0(y)=y$ for $y\ne y_0$. Being an isomorphism, $\tau_0$ is honest (Lemma~\ref{lem:retract}), so it extends to an endomorphism $\tau$ of $U=\ff Y$ by (F3). The fixed set $\Fix(\tau)=\{u:\tau(u)=u\}$ is a division subring of $U$ containing $Y\setminus\{y_0\}\supseteq s(Y)$, hence containing the division subring $D$ generated by $s(Y)$. Now $\sigma(U)\subseteq D$, because $\{u\in U:\sigma(u)\in D\}$ is a division subring containing $Y$. So $\tau$ is the identity on $\sigma(U)$, while $\tau(y_0)=y_0+1\ne y_0$. Hence $y_0\notin\sigma(U)$.
\end{proof}

\begin{remark}
Cohn uses precisely this set-up in \cite[\S5.9]{CohSF} and \cite[\S5]{Coh77}: a free algebra on generators indexed by $\Lambda\times\N^2$, an endomorphism shifting one index (which has a retraction and is therefore honest), a derivation shifting the other index, both extended to the free field. In \cite{Coh77} the non-surjectivity of the extended endomorphism is stated to be clear; Lemma~\ref{lem:nonsurj} supplies the argument.
\end{remark}

\section{Adjoining the solution of one equation}\label{sec:adjoin}

Throughout this section $K$ is a division ring with an endomorphism $\sigma$ and a $\sigma$-derivation $\delta$, and $C\in\M_n(K)$ is a fixed matrix. Let $A\supseteq K$ be a ring with an endomorphism and a $\sigma$-derivation extending those of $K$, which we again denote by $\sigma,\delta$; for row vectors and matrices over $A$ we apply $\sigma,\delta$ entrywise, and we put $T(v)=\sigma(v)C+\delta(v)$ for $v\in A^n$.

\begin{lemma}\label{lem:fixedstep}
Let $v\in A^n$ with $\sigma(v)=v$ and let $N\in\M_n(K)$. Then
\[
T(vN)=v\bigl(\sigma(N)C+\delta(N)\bigr)+\delta(v)N .
\]
\end{lemma}

\begin{proof}
Entrywise, $\sigma(vN)=\sigma(v)\sigma(N)=v\sigma(N)$ and $\delta(vN)=\sigma(v)\delta(N)+\delta(v)N=v\delta(N)+\delta(v)N$, by the Leibniz rule applied to each entry $\sum_l v_lN_{lj}$. Add.
\end{proof}

Define matrices $N_{i,j}\in\M_n(K)$ for $0\le j\le i$ by $N_{0,0}=I$ and
\[
N_{i+1,j}=\sigma(N_{i,j})C+\delta(N_{i,j})+N_{i,j-1}\qquad(0\le j\le i+1),
\]
with the conventions $N_{i,-1}=0$ and $N_{i,i+1}=0$. Thus $N_{i,i}=I$ for all $i$, and the $N_{i,j}$ depend only on $C$ and on $\sigma,\delta$ restricted to $K$.

\begin{lemma}\label{lem:iterate}
Let $m\ge1$ and let $x_0,\dots,x_{m-1}\in A^n$ satisfy $\sigma(x_j)=x_j$ for $0\le j\le m-1$ and $\delta(x_j)=x_{j+1}$ for $0\le j\le m-2$. Put $x_m=\delta(x_{m-1})$. Then for $0\le i\le m$,
\[
T^i(x_0)=\sum_{j=0}^{i}x_jN_{i,j}.
\]
\end{lemma}

\begin{proof}
Induction on $i$, the case $i=0$ being trivial. For $i<m$, every $x_j$ with $j\le i$ is fixed by $\sigma$ and satisfies $\delta(x_j)=x_{j+1}$ (for $j=m-1$ by definition of $x_m$), so Lemma~\ref{lem:fixedstep} gives
\[
T^{i+1}(x_0)=\sum_{j\le i}T(x_jN_{i,j})=\sum_{j\le i}x_j\bigl(\sigma(N_{i,j})C+\delta(N_{i,j})\bigr)+\sum_{j\le i}x_{j+1}N_{i,j}=\sum_{j\le i+1}x_jN_{i+1,j}.\qedhere
\]
\end{proof}

\begin{proposition}\label{prop:adjoin}
Let $Y$ be a set, $s\colon Y\to Y$ injective, $d\colon Y\to\ff Y$ a map, and let $K=\ff Y$ with the pair $(\sigma,\delta)$ of Proposition~\ref{prop:extend}. Let $f\in K[t;\sigma,\delta]$ be monic of degree $n\ge1$, let $C=C_f$, let $g=\sum_{i=0}^mg_it^i$ with $m\ge1$ and $g_m\ne0$, and let $e\in K^n$. Let $X=\{x_{l,j}:1\le l\le n,\ 0\le j\le m-1\}$ be a set of $nm$ new symbols, $Y'=Y\sqcup X$ and $K'=\ff{Y'}\supseteq K$. Write $x_j=(x_{1,j},\dots,x_{n,j})\in K'^n$ and define, in $K'^n$,
\[
\Phi=\sum_{i=0}^{m}g_i\sum_{j=0}^{\min(i,m-1)}x_jN_{i,j},\qquad \zeta=g_m^{-1}(e-\Phi),
\]
where the $N_{i,j}\in\M_n(K)$ are the matrices defined above. Extend $s$ to $s'\colon Y'\to Y'$ by $s'(x)=x$ for $x\in X$, and $d$ to $d'\colon Y'\to K'$ by
\[
d'(x_{l,j})=x_{l,j+1}\ (0\le j\le m-2),\qquad d'(x_{l,m-1})=\zeta_l ,
\]
and let $(\sigma',\delta')$ be the pair on $K'$ given by Proposition~\ref{prop:extend} for $(Y',s',d')$. Then $\sigma',\delta'$ restrict to $\sigma,\delta$ on $K$, and, computing $T_f$ in $K'$ with $\sigma',\delta'$,
\[
g(T_f)(x_0)=e .
\]
\end{proposition}

\begin{proof}
$Y\subseteq Y'$ is stable under $s'$ and $d'(Y)=d(Y)\subseteq\ff Y$, so the restriction statement is the last part of Proposition~\ref{prop:extend}. Hence $K'$ is a ring as in the standing assumptions of this section, with $A=K'$, and the matrices $N_{i,j}$ computed in $K$ are the ones appearing in Lemma~\ref{lem:iterate}. The vectors $x_0,\dots,x_{m-1}$ satisfy the hypotheses of Lemma~\ref{lem:iterate}, since $s'$ fixes $X$ and $d'(x_{l,j})=x_{l,j+1}$ for $j\le m-2$; and $x_m=\delta'(x_{m-1})=\zeta$. Therefore
\[
g(T_f)(x_0)=\sum_{i=0}^{m}g_i\sum_{j=0}^{i}x_jN_{i,j}=g_m\,x_mN_{m,m}+\Phi=g_m\zeta+\Phi=e .\qedhere
\]
\end{proof}

Note that $\zeta$ is an explicit element of $K'^n$ defined before $\sigma'$ and $\delta'$ are introduced; there is no circularity. The only role of the generators $x_{l,j}$ with $j\ge1$ is to serve as the values of $\delta'$ on the previous ones, so that the equation $g(T_f)(x_0)=e$ involves exactly one unknown quantity, $\delta'(x_{m-1})$, and that quantity is solved for.

\section{The construction}\label{sec:construction}

\begin{construction}\label{constr}
Let $k$ be a commutative field. We define recursively sets $Y_i$ with injective maps $s_i\colon Y_i\to Y_i$ and maps $d_i\colon Y_i\to\ff{Y_i}$, such that $Y_i\subseteq Y_{i+1}$, $s_{i+1}|_{Y_i}=s_i$ and $d_{i+1}|_{Y_i}=d_i$. We write $K_i=\ff{Y_i}$ with the pair $(\sigma_i,\delta_i)$ given by Proposition~\ref{prop:extend}, and $R_i=K_i[t;\sigma_i,\delta_i]$.
\begin{itemize}
\item $Y_0=\{y_j:j\in\N\}$, $s_0(y_j)=y_{j+1}$, $d_0=0$.
\item Given $(Y_i,s_i,d_i)$, let $\T_i$ be the set of all triples $(f,g,e)$ with $f\in R_i$ monic of degree $n\ge1$, $g\in R_i$ of degree $\ge1$ and $e\in K_i^n$. For each $\tau=(f,g,e)\in\T_i$, with $n=\deg f$ and $m=\deg g$, choose a set $X_\tau$ of $nm$ new symbols $x^\tau_{l,j}$, the sets $X_\tau$ being pairwise disjoint and disjoint from $Y_i$. Put $Y_{i+1}=Y_i\sqcup\bigsqcup_{\tau\in\T_i}X_\tau$. Let $s_{i+1}$ be $s_i$ on $Y_i$ and the identity on each $X_\tau$. Let $d_{i+1}$ be $d_i$ on $Y_i$, and on $X_\tau$ let it be the map $d'$ of Proposition~\ref{prop:adjoin} for the data $(Y_i,s_i,d_i)$ and $(f,g,e)$; its values lie in $\ff{Y_i\sqcup X_\tau}\subseteq\ff{Y_{i+1}}$.
\end{itemize}
Finally let $Y=\bigcup_iY_i$, $s=\bigcup_is_i\colon Y\to Y$, $d=\bigcup_id_i\colon Y\to\ff Y$ (using $\ff{Y_i}\subseteq\ff Y$), and let $\Kt=\ff Y$ with the pair $(\st,\dt)$ of Proposition~\ref{prop:extend}. The map $s$ is injective because each $s_i$ is injective and any two elements of $Y$ lie in a common $Y_i$.
\end{construction}

\begin{proof}[Proof of Theorem~\ref{thm:target}]
We show that $\Kt,\st,\dt$ have the required properties.

\emph{Compatibility.} For each $i$, the subset $Y_i\subseteq Y$ is stable under $s$ and $d(Y_i)\subseteq\ff{Y_i}=K_i$, so by Proposition~\ref{prop:extend} the subring $K_i\subseteq\Kt$ is stable under $\st,\dt$, and $(\st,\dt)$ restricts to $(\sigma_i,\delta_i)$ on $K_i$. Likewise $(\sigma_{i+1},\delta_{i+1})$ restricts to $(\sigma_i,\delta_i)$ on $K_i\subseteq K_{i+1}$, and for $\tau\in\T_i$ the pair $(\sigma_{i+1},\delta_{i+1})$ restricts on $\ff{Y_i\sqcup X_\tau}$ to the pair $(\sigma',\delta')$ of Proposition~\ref{prop:adjoin}, since $Y_i\sqcup X_\tau$ is stable under $s_{i+1}$ and $d_{i+1}$ maps it into $\ff{Y_i\sqcup X_\tau}$. Consequently $R_i$ is a subring of $\Rt=\Kt[t;\st,\dt]$, and for $f\in R_i$ monic the operator $T_f$ on $K_i^n$ is the restriction of the operator $T_f$ on $\Kt^n$. By Lemma~\ref{lem:subset}, $\Kt=\bigcup_iK_i$.

\emph{Solvability.} Let $f\in\Rt$ be monic of degree $n\ge1$, $g\in\Rt$ of degree $m\ge1$ and $e\in\Kt^n$. The finitely many coefficients of $f$, $g$ and $e$ lie in some $K_i$, so $\tau=(f,g,e)\in\T_i$. By Proposition~\ref{prop:adjoin}, applied inside $\ff{Y_i\sqcup X_\tau}\subseteq K_{i+1}\subseteq\Kt$, the vector $v=x^\tau_0=(x^\tau_{1,0},\dots,x^\tau_{n,0})$ satisfies $g(T_f)(v)=e$, and by compatibility this identity holds in $\Kt^n$.

\emph{Non-surjectivity.} $y_0\notin s(Y)$, since $s(Y_0)=\{y_j:j\ge1\}$ and $s$ fixes every generator outside $Y_0$. By Lemma~\ref{lem:nonsurj}, $y_0\notin\st(\Kt)$.

\emph{Countability.} Suppose $k$ is countable. $Y_0$ is countable. If $Y_i$ is countable then so are $K_i$ (Lemma~\ref{lem:generated}), $R_i$, $K_i^n$ and $\T_i$, and each $X_\tau$ is finite, so $Y_{i+1}$ is countable. Hence $Y$ and $\Kt$ are countable.
\end{proof}

\begin{remark}[a variant with a fixed set of symbols]\label{rem:fixed}
It is convenient, for instance when the construction is to be carried out formally, to fix the symbols in advance instead of creating the sets $X_\tau$ stage by stage. Take $Y=Y_0\sqcup(\N\times\N)$ with $s$ the shift on $Y_0$ and the identity on $\N\times\N$, let $U=\ff Y$, and let $Y_i=Y_0\sqcup(\{0,\dots,i-1\}\times\N)$. Define maps $d_i\colon Y\to U$ recursively, with $d_i=0$ outside $Y_i$ and $d_{i+1}=d_i$ on $Y_i$: given $d_i$, let $(\sigma_i,\delta_i)$ be the pair on all of $U$ determined by $(s,d_i)$ (Proposition~\ref{prop:extend}), let $K_i$ be the division subring generated by $Y_i$, which is stable under $\sigma_i,\delta_i$ by the restriction clause, and enumerate the (countably many) triples $(f,g,e)$ over $K_i$ by natural numbers $c$; the symbols $(i,\langle c,l,j\rangle)$, with $\langle\cdot\rangle$ a pairing function, play the role of $x^\tau_{l,j}$, and $d_{i+1}$ is defined on them by the recipe of Proposition~\ref{prop:adjoin}, computed with $\sigma_i,\delta_i$; on symbols $(i,c')$ not of this form put $d_{i+1}(i,c')=0$. Finally $d=\bigcup_id_i$ and $(\st,\dt)$ is the pair determined by $(s,d)$. Since $d_j$ and $d$ agree on $Y_i$ for $j\ge i$, the pairs $(\sigma_j,\delta_j)$ and $(\st,\dt)$ all agree on $K_i$ by the uniqueness part of Proposition~\ref{prop:extend}, and the matrices $N_{i,j}$ of Proposition~\ref{prop:adjoin}, whose entries lie in $K_i$, are the same whether computed with $(\sigma_i,\delta_i)$ or with $(\st,\dt)$. The proof of Theorem~\ref{thm:target} then goes through verbatim; the unused symbols are extra free generators fixed by $\st$ and killed by $\dt$, and do no harm. This variant avoids the identifications of Lemma~\ref{lem:subset} between the fields $\ff{Y_i}$, at the cost of an enumeration.
\end{remark}

\begin{proof}[Proof of the Main Theorem]
Let $k$ be countable and let $\Kt,\st,\dt$ be as in Theorem~\ref{thm:target}; then $\Rt=\Kt[t;\st,\dt]$ is countable. By Proposition~\ref{prop:criterion}, $\Rt$ is a simple principal left ideal domain and a left PCI domain, hence a left V-domain (by Lemma~\ref{lem:pciv}, or by \cite[Theorem~3.2]{JLL09}); in particular it is left Noetherian and left hereditary. By Proposition~\ref{prop:notore}, with $c=y_0\notin\st(\Kt)$, $\Rt$ is not right Ore and $t\Rt\cap y_0t\Rt=0$; so by Lemma~\ref{lem:countable} no simple right $\Rt$-module is injective, $\Rt$ is not a right V-ring, and by Lemma~\ref{lem:pciv} ($\Rt$ is a domain and not a division ring, as $t$ is not a unit) it is not a right PCI ring. (That $\Rt$ is not right PCI also follows, without countability, from Proposition~\ref{prop:notore} and \cite[6.17]{CF75}.) The statements about $\Rt^{\op}$ are the mirror images; that a right PCI domain is simple, right Noetherian and right hereditary is \cite{Fai73} and \cite[Theorem~1]{Dam79}.
\end{proof}

\begin{corollary}\label{cor:answers}
Let $\Rt$ be as in the Main Theorem.
\begin{enumerate}
\item A right PCI ring need not be left PCI \textup{(}\cite{Fai73}, \cite[\S7, (10)]{CF75}, \cite{Dam79}\textup{)}: $\Rt^{\op}$.
\item A right PCI domain need not be left Ore, and a right V-domain need not be left Ore \textup{(}\cite[\S7, (7)]{CF75}\textup{)}: $\Rt^{\op}$.
\item A left V-domain need not be a right V-domain \textup{(}\cite[\S7, (10)]{CF75}, \cite[\S1]{JLL09}, \cite[Remark~2.8]{BDV18}\textup{)}: $\Rt$.
\item There is an Ore extension $K[t;\sigma,\delta]$ which is a left V-domain with $\sigma$ not onto \textup{(}\cite[\S6]{JLL09}\textup{)}: $\Rt$.
\end{enumerate}
\end{corollary}

\section{Remarks}

\begin{remark}[on the shape of the construction]
Three features of the construction are forced.
\begin{enumerate}
\item The derivation cannot be omitted or made inner. If $\delta=0$ then $Rt$ is a proper two-sided ideal of $K[t;\sigma]$ (as $tR\subseteq Rt$), so the ring is not simple. If $\delta$ is inner, $\delta(a)=ca-\sigma(a)c$, the substitution $t'=t-c$ gives $K[t;\sigma,\delta]=K[t';\sigma]$. In particular $K$ cannot be commutative when $\sigma\ne\id$, since then every $\sigma$-derivation is inner \cite[\S2]{Coh77}.
\item The adjoined solutions must be free over $K$. If a solution $x$ of $g(T_f)(x)=e_1$ satisfied a relation $xa=\omega(a)x$ ($a\in K$) with $\omega$ an endomorphism of $K$, applying $\delta$ and comparing coefficients would force $\sigma\circ\omega$ to be inner, which is impossible for non-surjective $\sigma$. This is why the extension is a free field rather than an Ore or a crossed-product extension.
\item The closure has to be iterated: the data $(f,g)$ range over polynomials whose coefficients are elements of the field being built, and the values $\zeta$ of the derivation on the new generators involve $\sigma$ and $\delta$ applied to those coefficients.
\end{enumerate}
By contrast, the following are matters of convenience: taking the new generators to be fixed by $\sigma$ (which makes the computation of $g(T_f)$ a two-line induction and requires only $nm$ generators per equation), solving for every right-hand side $e$ rather than only for $e_1$ (which is all the criterion $R=Rf+gR$ of \cite[Theorem~3.2]{JLL09} asks for, while every $e$ is what the Baer-criterion argument after Proposition~\ref{prop:criterion} uses), and choosing $\ff{Y_0}$ with the shift as the starting field.
\end{remark}

\begin{remark}[on the literature used]
Apart from the results on Ore extensions and PCI rings quoted in Section~\ref{sec:ore}, the proof uses from Cohn's theory only Facts~\ref{facts}: free algebras are firs, semifirs have universal fields of fractions with the universal property of a localization, and honest maps extend to them. Lemma~\ref{lem:generated}, Lemma~\ref{lem:retract} and Proposition~\ref{prop:extend} are proved from these. The extension of a derivation to a universal field of fractions (Proposition~\ref{prop:extend}, Steps~2--3) is the argument of \cite[Theorem~7.5.17]{CohFIR}, with a twist by $\sigma$ inserted. No coproduct theorems are needed.
\end{remark}

\begin{remark}[further properties of $\Rt$]
By \cite[Theorem~6.2]{JLL09}, the maximal right quotient ring of the left V-domain $\Rt$ is directly infinite, and by \cite[Proposition~6.1]{JLL09} no cyclic right module $\Rt/f\Rt$ with $f\ne0$ is injective and every principal right ideal of $\Rt$ is closed. The proof of Lemma~\ref{lem:countable} shows more: every non-zero injective right $\Rt$-module $E$ satisfies $|E|^{\aleph_0}\le|E|$, hence $|E|\ge2^{\aleph_0}$. If $k$ is uncountable one still obtains, by Theorem~\ref{thm:target}, a left PCI domain $\Kt[t;\st,\dt]$ which is not right Ore and hence not right PCI; only the failure of the right V-property was derived from countability.
\end{remark}

\begin{thebibliography}{JLL09}

\bibitem[BDV18]{BDV18} M.~Behboodi, A.~Daneshvar, M.~R.~Vedadi, \emph{Several generalizations of the Wedderburn--Artin theorem with applications}, Algebr. Represent. Theory \textbf{21} (2018), 1333--1342.

\bibitem[Boy74]{Boy74} A.~K.~Boyle, \emph{Hereditary QI-rings}, Trans. Amer. Math. Soc. \textbf{192} (1974), 115--120.

\bibitem[BG75]{BG75} A.~K.~Boyle, K.~R.~Goodearl, \emph{Rings over which certain modules are injective}, Pacific J. Math. \textbf{58} (1975), 43--53.

\bibitem[Coh77]{Coh77} P.~M.~Cohn, \emph{A construction of simple principal right ideal domains}, Proc. Amer. Math. Soc. \textbf{66} (1977), 217--222.

\bibitem[CohSF]{CohSF} P.~M.~Cohn, \emph{Skew Fields: Theory of General Division Rings}, Encyclopedia of Mathematics and its Applications \textbf{57}, Cambridge University Press, 1995.

\bibitem[CohFIR]{CohFIR} P.~M.~Cohn, \emph{Free Ideal Rings and Localization in General Rings}, New Mathematical Monographs \textbf{3}, Cambridge University Press, 2006.

\bibitem[Coz70]{Coz70} J.~H.~Cozzens, \emph{Homological properties of the ring of differential polynomials}, Bull. Amer. Math. Soc. \textbf{76} (1970), 75--79.

\bibitem[CF75]{CF75} J.~H.~Cozzens, C.~Faith, \emph{Simple Noetherian Rings}, Cambridge Tracts in Mathematics \textbf{69}, Cambridge University Press, 1975.

\bibitem[Dam79]{Dam79} R.~F.~Damiano, \emph{A right PCI ring is right Noetherian}, Proc. Amer. Math. Soc. \textbf{77} (1979), 11--14.

\bibitem[Fai73]{Fai73} C.~Faith, \emph{When are proper cyclics injective?}, Pacific J. Math. \textbf{45} (1973), 97--112.

\bibitem[JLL09]{JLL09} S.~K.~Jain, T.~Y.~Lam, A.~Leroy, \emph{Ore extensions and V-domains}, in: Rings, Modules and Representations, Contemp. Math. \textbf{480}, Amer. Math. Soc., 2009, 249--262.

\bibitem[JS09]{JS09} S.~K.~Jain, A.~K.~Srivastava, \emph{Rings determined by the properties of cyclic modules and duals: a survey-I}, preprint, 2009, \url{https://leroy.perso.math.cnrs.fr/Talks/Jain.pdf}.

\bibitem[Law77]{Law77} J.~Lawrence, \emph{A countable self-injective ring is quasi-Frobenius}, Proc. Amer. Math. Soc. \textbf{65} (1977), 217--220.

\end{thebibliography}

\end{document}
